\documentclass[12pt]{report}

% package to activate greek language 
\usepackage[greek,english]{babel}
\usepackage{indentfirst}

% DOCUMENT LAYOUT
\usepackage{geometry} 
\geometry{a4paper, textwidth=5.5in, textheight=8.5in, marginparsep=7pt, marginparwidth=.6in}
\setlength\parindent{8mm}
\setlength\parskip{5mm}

% FONTS
\usepackage{xltxtra}
\usepackage[table]{xcolor}
\defaultfontfeatures{Mapping=tex-text}
\usepackage{libertine}

% package to handle graphics
\usepackage{graphicx}
% package to handle multiple figures in a minipage
\usepackage{subfigure}
%package to activate XeTeX font manager
\usepackage{fontspec}

% PDF SETUP
% ---- FILL IN HERE THE DOC TITLE AND AUTHOR
\usepackage[driverfallback=dvipdfmx, bookmarks, colorlinks, breaklinks, pdftitle={Efi Kaltirimidou},pdfauthor={efikalti}]{hyperref}  
\hypersetup{linkcolor=blue,citecolor=blue,filecolor=black,urlcolor=blue} 

% change captions especially for greek language - if the document is in ENGLISH, they should vanish
\addto\captionsenglish{%
  \renewcommand\prefacename{Πρόλογος}%
  \renewcommand\refname{Αναφορές}%
  \renewcommand\abstractname{Περίληψη}%
  \renewcommand\bibname{Βιβλιογραφία}%
  \renewcommand\chaptername{Κεφάλαιο}%
  \renewcommand\appendixname{Παράρτημα}%
  \renewcommand\contentsname{Περιεχόμενα}%
  \renewcommand\listfigurename{Κατάλογος Σχημάτων}%
  \renewcommand\listtablename{Κατάλογος Πινάκων}%
  \renewcommand\indexname{Ευρετήριο}%
  \renewcommand\figurename{Σχήμα}%
  \renewcommand\tablename{Πίνακας}%
  \renewcommand\partname{Μέρος}%
  \renewcommand\enclname{Συνημμένα}%
  \renewcommand\ccname{Κοινοποίηση}%
  \renewcommand\headtoname{Προς}%
  \renewcommand\pagename{Σελίδα}%
  \renewcommand\seename{βλέπε}%
  \renewcommand\alsoname{βλέπε επίσης}%
  \renewcommand\proofname{Απόδειξη}%
  }

%%%%%%%%% END OF PREAMBLE %%%%%%%%%%%%

\begin{document}

% make title out of \author, \title, \date  specified in the preamble
\begin{titlepage}
%
%
\begin{figure}[ht]
\begin{minipage}{0.15\textwidth}
\vspace{-10mm}
\includegraphics[width=20mm]{../images/aristotle_logo.png}
\end{minipage}
%
\begin{minipage}{0.70\textwidth}
{
ΑΡΙΣΤΟΤΕΛΕΙΟ ΠΑΝΕΠΙΣΤΗΜΙΟ\\[0.5\baselineskip]
ΘΕΣΣΑΛΟΝΙΚΗΣ\\[0.5\baselineskip]
Τμήμα Πληροφορικής
}
\end{minipage}
%
\end{figure}

\vspace{10mm}
\hbox{ % Horizontal box
\hspace*{0.08\textwidth} % Whitespace to the left of the title page
\rule{1.5pt}{0.80\textheight} % Vertical line
\hspace*{0.02\textwidth} % Whitespace between the vertical line and title page text
\parbox[b]{0.90\textwidth}{ % Paragraph box which restricts text to less than the width of the page


{\large \textsc{Πτυχιακή Εργασία}}\\[\baselineskip] % Tagline or further description
{\noindent\LARGE\bfseries
Σχεδίαση και Υλοποίηση Web
Πλατφόρμας για την Υποβολή Επιστημονικών Εφαρμογών σε Κατανεμημένα Συστήματα
}\\[2\baselineskip] % Title

{\Large Ευφροσύνη Καλτιριμίδου}\\[3\baselineskip] % Author name

{\em Επιβλέπων:}\\[0.25\baselineskip]
{\large Αθηνά Βακάλη}\\[0.25\baselineskip]
{\large Καθηγήτρια}

\vspace{0.15\textheight} % Whitespace between the title block and the publisher
{\noindent Θεσσαλονίκη 2017}\\[\baselineskip] % Publisher and logo
}}
%
%
\end{titlepage}

%\newpage
\thispagestyle{empty}
\phantom{~}
\vfill


% Create table of contents
\tableofcontents
%\newpage
\thispagestyle{empty}
\phantom{~}
\vfill

% Intro
\section{Περίληψη}


% Frontend
% print no page number
\thispagestyle{empty}

\chapter{Front - End Architecture}

To frontend είναι το πιο σημαντικό κομμάτι της εφαρμογής για τον χρήστη. Αν η εφαρμογή ήταν ένα σπίτι, το frontend θα ήταν η εξωτερική πρόσοψη του σπιτού και η οικοδεσπότης, ο οποίος καλωσορίζει τους ανθρώπους μέσα.  Το frontend είναι το κομμάτι της εφαρμογής που χρησιμοποιεί ο χρήστης για να αποκτήσει, να μεταβάλλει και να αποθηκεύσει δεδομένα. Στην συγκεκριμένη εργασία, το frontend είναι αποκλειστικά μια web εφαρμογή.
\newline
Όλη η επίδραση με τον χρήστη γίνετε μέσω του frontend, γιαυτό η σχεδίαση του αποτελεί το δεύτερο βήμα στην υλοποίηση της εφαρμογής. Η τεχνική υλοποίηση του θα γίνει μετά την τεχνική υλοποίηση του backend καθώς η λειτουργικότητα του στηρίζετε στα API calls που προσφέρει το backend.

% leave 60mm empty space below
\vspace{60mm}

\section{Στόχοι της σχεδίασης}
Οι αρχές που ακολουθεί η σχεδίαση του frontend χωρίζετε σε δύο ειδών, τις γενικές αρχές και τις αρχές ασφάλειας.
\begin{description}
\item [Γενικές αρχές] Αναφέρονται σε όλο το frontend και στοχεύουν να κάνουν την εφαρμογή απλή και εύκολη για τον χρήστη.
\newline
Μερικές γενικές αρχές:


\begin{itemize}
  \item Σε κάθε σελίδα του frontend, θα πρέπει να είναι ξεκάθαρες οι διαθέσιμες ενέργειες του χρήστη.
  \item Σε οποιαδήποτε σελίδα βρίσκετε ο χρήστης θα πρέπει να μπορεί να πάει σε οποιαδήποτε άλλη σελίδα της εφαρμογής με μία μόνο κίνηση.
   \item Για κάθε ενέργεια του χρήστη θα πρέπει να υπάρχει μια άμεση απάντηση της εφαρμογής, είτε με αλλαγή περιεχομένου, είτε με ενημερωτικό μύνημα κ.τ.λ.
   \item Σε περίπτωση αποτυχίας κάποιας ενέργειας, θα πρέπει να ενημερώνετε με λεπτομέρεια ο χρήστης ώστε να γνωρίζει πως να συνεχίσει. Αυτό περιλαμβάνει την χρήση μυνημάτων που θα εμφανίζονται στην οθόνη, το focus της σελίδας σε συγκεκριμένα πεδία ή δεδομένα κ.τ.λ.
\end{itemize}


\item [Αρχές ασφάλειας] Οι αρχές ασφάλειας αναφέρονται σε συγκεκριμένα κομμάτια του frontend και είναι κυρίως κομμάτι της ασφάλειας των δεδομένων.

\begin{itemize}
  \item Η εφαρμογή θα χρησιμοποιεί προτόκολλα ασφάλειας για την μεταφορά των δεδομένων του χρήστη.
  \item Ο χρήστης θα παραμένει συνδεμένος για πεπερασμένο χρονικό διάστημα και στη συνέχεια θα πρέπει να ανανεώσει την σύνδεση του.
\end{itemize}

\end{description}

\section{Τεχνικά χαρακτηριστικά}

Οι τεχνολογίες που χρησιμοποιούνται σε κάθε εφαρμογή ποικίλουν ανάλογα με τα άτομα που την υλοποιήσανε, τα χαρακτηριστικά και τις δυνατότητες της εφαρμογής.
\newline
Για τις web εφαρμογές, δύο τεχνολογίες αποτελούν αναπόσπαστα κομμάτια τους:

\begin{description}
\item [Hypertext Markup Language (\textbf{HTML})] Αποτελεί την γλώσσα σήμανσης που δημιουργεί τις σελίδες της εφαρμογής.

\item [Cascading Style Sheets (\textbf{CSS})] Η οποία περιγράφει την μορφοποίηση του κάθε αντικειμένου στη σελίδα. 
\end{description}


Σε συνδυασμό με μία client-based γλώσσα, όπως η Javascript, αποτελούν τα συστατικά στοιχεία κάθε web εφαρμογή σήμερα.

\begin{description}
\item [JavaScript (\textbf{JS})] Είναι μια client-based γλώσσα προγραμματισμού, η οποία δίνει δυνατότητα για δυναμικές ιστοσελίδες και εφαρμογές στο διαδίκτυο.
\end{description}

Για την υλοποίηση της εφαρμογής χρησιμοποιήθηκε το web framework \textbf{AngularJS}, ένα δομικό framework για δυναμικές web εφαρμογές. Χρησιμοποιεί την HTML ως γλώσσα προτύπου και επεκτείνει τη σύνταξη της για να εκφράσει σαφώς και συνοπτικά τα στοιχεία της εφαρμογής. Διάφορες  λειτουργίες της AngularJS μειώνουν ένα μεγάλο μέρος του κώδικα που θα έπρεπε να γραφτεί. Βασίζετε στη Javascript, και την επεκτείνει για να προσθέσει επιπλεόν λειτουργίες προς διευκόλυνση του προγραμματιστή.

\subsection{Εκδόσεις}

Η εφαρμογή έχει υλοποιηθεί και δοκιμαστεί, στο κομμάτι του frontend με τις παρακάτων εκδόσεις:

\begin{description}
\item [HTML] 5
\item [CSS] 3
\item [JavaScript] 5
\item [AngularJS] 1.6.*
\end{description}


% leave 60mm empty space below
\vspace{60mm}

\section{Λειτουργίες Frontent}

Στο frontend της εφαρμογής έχουν πρόσβαση όλοι οι χρήστες της  \textit{Ιδρυματικής Συστοιχίας ΑΠΘ} με τα στοιχεία πιστοποίησης του λογαριασμού τους στο \textit{GRID}. Έχουν πρόσβαση στις παρακάτω λειτουργίες:

\begin{description}
  \item[$\bullet$ Σύνδεση] O χρήστης θα παρέχει το username και password που έχει πάρει κατά την εγγραφή του ως μέλος της \textit{Ιδρυματικής Συστοιχίας ΑΠΘ}.\textit{Προσοχή} , η εφαρμογή δεν παρέχει μέθοδο εγγραφής. Η εφαρμογή θα στείλει τα δεδομένα με ασφαλή προτόκολλα στο backend για επιβεβαίωση. Αν τα δεδομένα δεν είναι σωστά θα εμφανίσει το ανάλογο μύνημα λάθους αλλιώς η εφαρμογή θα μεταβεί στο dashboard του χρήστη.
  \item[$\bullet$ Έλεγχος εργασιών] Ο χρήστης μπορεί να δει τις εργασίες που έχει κάνει submit στον cluster, σε τι κατάσταση βρίσκονται και να τις φιλτράρει βάση κατάστασης, ημερομηνίας υποβολής και ονόματος.
    \item[$\bullet$ Υποβολή εργασίας] Σε αυτή την λειτουργία ο χρήστης μπορεί να επιλέξει να υποβάλει το εκτελέσιμο της εργασίας από τον τοπικό του δίσκο ή να δημιουργήσει η εφαρμογή ένα εκτελέσιμο για αυτόν με βάση παραμέτρους που θα επιλέξει. Επίσης μπορεί να υποβάλει αρχεία εισόδου από τον τοπικό του δίσκο ή να δώσει το path που βρίσκονται μέσα στη συστοιχία.
    \item[$\bullet$ Παραλαβή αποτελεσμάτων] Αν κάποια εργασία έχει εκτελεστεί επιτυχώς, μπορεί να επιλέξει να κατεβάσει τα αποτελέσμα της στον τοπικό του δίσκο μέσω της εφαρμογής.
\end{description}








%\newpage
\thispagestyle{empty}
\phantom{~}
\vfill


% API
% print no page number
\thispagestyle{empty}

\chapter{Service layer - Restful API}

\par
Τα κομμάτια του frontend και του backend, παρόλο που είναι ανεξάρτητα το ένα από το άλλο, για την σωστή λειτουργία της εφαρμογής είναι απαραίτητο να επικοινωνήσουν μεταξύ τους. Την επικοινωνία τους αναλαμβάνει ένα τρίτο κομμάτι της εφαρμογής που ονομάζετε \textbf{Service Layer}.\par Σε πολλές εφαρμογές, το \textbf{Service Layer} υλοποιήτε από το κομμάτι του backend καθώς οι ίδιες μέθοδοι που κάνουν την επικοινωνία backend με frontend κάνουν και την επικοινωνία backend με τη βάση δεδομένων.Και σε αυτήν την περίπτωση το \textbf{Service Layer} έχει υλοποιηθεί μαζί με το backend.\par Σε γενικές γραμμές, πρόκειται για ένα σύνολο σαφώς καθορισμένων μεθόδων επικοινωνίας το οποίο για συντομία το λέμε \textbf{API (application programming interface)}.Το API ορίζει τον σωστό τρόπο για έναν προγραμματιστή να γράψει ένα πρόγραμμα που ζητά υπηρεσίες από λειτουργικό σύστημα (OS) ή άλλη εφαρμογή. Τα API υλοποιούνται με λειτουργικές κλήσεις που αποτελούνται από ρήματα και ουσιαστικά. Η απαιτούμενη σύνταξη περιγράφεται συνήθως στο documentation της κάθε εφαρμογής.


% leave 60mm empty space below
\vspace{60mm}

\section{Περιγραφή}
	Τα API αποτελούνται από δύο στοιχεία. Το πρώτο είναι ένα σύνολο κανόνων που περιγράφει τον τρόπο ανταλλαγής πληροφοριών μεταξύ προγραμμάτων, που γίνεται με τη μορφή αίτησης για επεξεργασία και επιστροφής των απαραίτητων δεδομένων. Το δεύτερο είναι το λογισμικού γραμμένο με τις προδιαγραφές που ορίστηκαν από τους κανόνες και δημοσιεύεται για χρήση.

Το λογισμικό που επιθυμεί να αποκτήσει πρόσβαση στις δυνατότητες και τις δυνατότητες του API λέγεται ότι καλεί το API και το λογισμικό που υλοποιεί το API λέγεται ότι το δημοσιεύει.

\section{Σχεδίαση του API}
	Ο σωστός σχεδιασμός του API βασίζετε στις αρχές της διαφάνειας και της ασφάλειας των δεδομένων. Το API θα πρέπει να αποκρύπτει τις μεθόδους και την πολυπλοκότητα της επικοινωνίας από και προς τις δύο υπηρεσίες.Έτσι, ο σχεδιασμός του API επιχειρεί να παρέχει μόνο τα εργαλεία που θα περίμενε να παραλάβει το frontend.
\newline

\section{Είδη APIs}
Υπάρχουν τριών ειδών APIs. 

\begin{description}
\item[Local APIs] Eίναι η αρχική μορφή, από την οποία ήρθε το όνομα. Προσφέρουν υπηρεσίες OS ή middleware σε προγράμματα εφαρμογών. Το API της Microsoft .NET, το TAPI (API τηλεφωνίας) για φωνητικές εφαρμογές και τα API πρόσβασης σε βάσεις δεδομένων είναι  Local API. 
\item[Web APIs] Αυτά τα API έχουν σχεδιαστεί για να αντιπροσωπεύουν ευρέως χρησιμοποιούμενους πόρους όπως οι σελίδες HTML και έχουν πρόσβαση μέσω ενός απλού πρωτοκόλλου HTTP. Οποιαδήποτε διαδικτυακή διεύθυνση URL καλεί ένα web API. Τα API του Web ονομάζονται συχνά REST  ή RESTful επειδή ο εκδότης των διασυνδέσεων REST δεν αποθηκεύει εσωτερικά δεδομένα μεταξύ αιτήσεων. Ως εκ τούτου, τα αιτήματα πολλών χρηστών μπορούν να αναμειχθούν όπως και στο διαδίκτυο.
\item[Program APIs] Τα API προγραμμάτων βασίζονται σε τεχνολογία κλήσης εξ αποστάσεως (Remote Procedure Call) (RPC), που καθιστά ένα εξ αποστάσεως συστατικό του προγράμματος να φαίνεται τοπικό για το υπόλοιπο λογισμικό. Για παράδειγμα η σειρά API της Microsoft WS.
\end{description}

Καθώς η συγκεκριμένη εφαρμογή έχει επικοινωνία με web application έχει υλοποιηθεί ένα web API.

\section{REST - RESTful API}
Οι REST (Representational state transfer) ή RESTful υπηρεσίες, όπως η συγκεκριμένη, αποτελούν έναν τρόπο παροχής λειτουργικότητας μεταξύ συστημάτων στο Διαδίκτυο. Οι υπηρεσίες Web που είναι συμβατές με τους κανόνες που ορίζει το REST επιτρέπουν στα συστήματα να έχουν πρόσβαση και να χειρίζονται τα δεδομένα του διαδικτύου χρησιμοποιώντας ένα ομοιόμορφο και προκαθορισμένο σύνολο λειτουργιών. Υπάρχουν και άλλες μορφές υπηρεσιών με τους δικούς τους κανόνες και λειτουργίες, όπως το WSDL και το SOAP.
\par Σε μια RESTful υπηρεσία , τα αιτήματα \textit{(requests)} που υποβάλλονται στο URL ενός πόρου θα προκαλέσουν μια απάντηση \textit{(response)} που μπορεί να είναι XML, HTML, JSON ή κάποια άλλη καθορισμένη μορφή. Η απόκριση μπορεί να επιβεβαιώσει ότι έχουν γίνει ορισμένες αλλαγές στον αποθηκευμένο πόρο και μπορεί να παρέχει συνδέσεις υπερκειμένου με άλλους σχετικούς πόρους ή συλλογές πόρων. Χρησιμοποιώντας HTTP, όπως είναι πιο συνηθισμένο, το είδος των διαθέσιμων λειτουργιών περιλαμβάνει εκείνες που είναι προκαθορισμένες από τα ρήματα HTTP GET, POST, PUT, DELETE κ.τ.λ

\subsection{Χαρακτηριστικά του REST}
Οι κανόνες που ορίζονται από το REST εξυπηρετούν στην διασφάλιση ορισμένων πλεονεκτημάτων που ορίζουν τις RESTful υπηρεσίες όταν υλοποιούνται σωστά.Τα παρακάτω είναι μερικά από τα χαρακτηριστικά που προσπαθεί να διασφαλίσει το REST
\begin{itemize}
\item Performance 
\item Scalability
\item Simplicity
\item Modifiability
\item Reliability
\end{itemize}


\subsection{Από URL σε HTTP}
Ένας από τους πιο βασικούς κανόνες που ορίζει το REST είναι η αντιστοίχηση των URL με συγκεκριμένα HTTP ρήματα και λειτουργίες.
Ο παρακάτω πίνακας περιγράφει ένα από τα αντικείμενα της εφαρμογής, το \textit{job} και πως αντιστοιχίζονται τα URL της web εφαρμογής με τα HTTP calls του API.

\begin{center}
\begin{tabular}{ |c|c|c|c| } 
\hline
URL & HTTP GET & col3 \\
\hline
http://{backend}/api/jobs/ & cell5 & cell6 \\ 
2 & cell8 & cell9 \\ 
\hline
\end{tabular}
\end{center}








%\newpage
\thispagestyle{empty}
\phantom{~}
\vfill

\end{document}